% Page 19 — Alban en sa « légende dorée »
% À intégrer dans le document principal (XeLaTeX).
% Nécessite le package lettrine.

\begin{center}
    {\Large\itshape Alban en sa « légende dorée »}
\end{center}

\medskip

\lettrine[lines=2]{M}{ais} devinez donc ce qui vient d’arriver ? Alain de Nougaret,
Baron d’Alban a signé le contrat qui le livre à Satan pour l’éternité.
Lucifer regarde avec un sourire féroce celui qui vient de se laisser prendre
à ses infernales propositions. Oui, c’est vrai, Alain a signé de son sang
le parchemin diabolique, mais Alain est libre. Si vous pensez peut-être
qu’il va profiter de cette liberté pour faire le bien et se racheter ?
Trois fois hélas ! c’est de nouveau une vie de pillage qui recommence.
Et de jour et de nuit, il est à l’affût de quelque victime. Satan a bien
fait les choses. Alain de Nougaret arpente les chemins du crime, dévore
les routes spacieuses qui mènent à l’enfer éternel. Mon Dieu, mon Dieu,
jusqu’où ira-t-il ? Et la Madone de pitié l’aurait-elle définitivement
abandonné ?

\medskip

\lettrine[lines=2]{N}{on}, la Vierge n’abandonne jamais, comme vous allez le voir,
mais la folie des hommes peut aller loin. Et c’est par une soirée de ripaille
et de débauche qu’Alain va être pris dans les filets d’une miséricorde infinie.
Il festoie en son château de Nougaret tandis que ses paysans à nouveau révoltés
assiègent son repaire et de toutes parts y mettent le feu. Femmes, enfants,
domestiques, tout le monde va périr dans les flammes et lui-même sera grillé.
Qui l’eût cru ? Alain est le seul à s’être échappé du château qui flambe et
le voilà meurtri mais bien vivant au milieu d’un fourré… Alors…
